% =====================================================================
%  CAPÍTULO 2 — FUNDAMENTAÇÃO TEÓRICA
% =====================================================================

\chapter{FUNDAMENTAÇÃO TEÓRICA}
\label{cap:fundamentacao}

Este capítulo apresenta os conceitos teóricos necessários à compreensão do
\argus{}, abrangendo desde os protocolos de comunicação sem fio de baixo
consumo até os modelos físicos e fisiológicos utilizados no cálculo de
telemetria, bem como os princípios de engenharia de software que guiam a
arquitetura do sistema.

\section{Bluetooth Low Energy e perfis GATT}
\label{sec:ble-gatt}

O \textit{Bluetooth Low Energy} (BLE) é uma tecnologia de comunicação sem fio
de curto alcance e baixo consumo energético, projetada para dispositivos que
necessitam transmitir pequenas quantidades de dados de forma intermitente,
como sensores esportivos \cite{bluetooth2017}. No contexto do ciclismo, sensores
de frequência cardíaca, medidores de potência e rolos de treino inteligentes
utilizam o BLE para transmitir telemetria em tempo real.

O \textit{Generic Attribute Profile} (GATT) define a estrutura de dados sobre a
qual os serviços são expostos. Cada serviço é identificado por um UUID e contém
características (\textit{characteristics}) que podem ser lidas, escritas ou
assinar notificações. Os serviços relevantes ao \argus{} são:

\begin{itemize}
    \item \textbf{Heart Rate Service} (\texttt{0x180D}): expõe a frequência
        cardíaca instantânea do atleta, geralmente em batimentos por minuto
        (BPM) \cite{bluetooth2017}.
    \item \textbf{Cycling Power Service} (\texttt{0x1818}): transmite a
        potência instantânea (em watts) e, opcionalmente, cadência e
        distribuição do torque aplicado aos pedais.
    \item \textbf{Fitness Machine Service} (FTMS, \texttt{0x1826}): permite o
        controle bidirecional de equipamentos de fitness, incluindo a leitura
        de telemetria do rolo e o envio de comandos de resistência, habilitando
        modos como ERG (potência alvo fixa) e simulação de inclinação
        \cite{bluetooth2017}.
\end{itemize}

Além do FTMS, rolos de treino de diferentes fabricantes podem expor o protocolo
FE-C (ANT+), que também é suportado pelo \argus{} para controle de resistência
em modos ERG e de altitude simulada.

\section{Protocolos e formatos de dados esportivos}
\label{sec:formatos}

\subsection{GPX (\textit{GPS Exchange Format})}
\label{sec:gpx}

O GPX é um formato aberto baseado em XML, criado pela Topografix, para o
intercâmbio de dados de posicionamento geográfico (latitude, longitude e
elevação) \cite{gpx2004}. No \argus{}, rotas GPX são importadas para extrair o
perfil de elevação e, a partir da diferença de altitude entre pontos
consecutivos, calcular o gradiente percentual da via, que alimenta o motor de
simulação física.

\subsection{FIT (\textit{Flexible and Interoperable Data Transfer})}
\label{sec:fit}

O formato FIT, especificado pela Garmin e pela aliança ANT+, é um formato
binário compacto e interoperável para armazenamento de atividades esportivas
\cite{garmin2024}. Arquivos FIT são amplamente aceitos por plataformas como
Strava, Garmin Connect e TrainingPeaks. O \argus{} gera arquivos FIT ao final
de cada sessão, permitindo que os treinos sejam importados por essas
plataformas sem conversão.

\subsection{ZWO (\textit{Zwift Workout})}
\label{sec:zwo}

O ZWO é o formato de plano de treino estruturado utilizado pela plataforma
Zwift, definindo intervalos com duração e alvo de potência (em percentual do
FTP) em modo ERG \cite{zwift2024}. O \argus{} importa e gera arquivos ZWO,
permitindo que planos de treino elaborados pelo assistente de IA sejam
carregados diretamente no simulador.

\section{Modelagem física do ciclismo: a equação de Martin}
\label{sec:equacao-martin}

O movimento de uma bicicleta pode ser modelado pelo balanço de forças que
atuam sobre o conjunto ciclista--bicicleta. O modelo proposto por Martin
\cite{martin1998} relaciona a potência aplicada pelo ciclista ($P$) à
velocidade ($v$), considerando as forças de arrasto aerodinâmico, resistência
ao rolamento, componente gravitacional do gradiente e aceleração linear,
conforme a Equação~\ref{eq:martin}:

\begin{equation}\label{eq:martin}
P = v \cdot \left( \frac{1}{2} \rho C_d A \, v^2 + m g C_{rr} + m g \sin\theta + m a \right) \cdot \eta^{-1}
\end{equation}

\noindent em que:
\begin{itemize}
    \item $\rho$ é a densidade do ar (\SI{1,225}{\kilo\gram\per\cubic\meter}
        ao nível do mar);
    \item $C_d A$ é o coeficiente de arrasto aerodinâmico multiplicado pela
        área frontal;
    \item $m$ é a massa total (ciclista + bicicleta);
    \item $g$ é a aceleração gravitacional
        (\SI{9,81}{\meter\per\square\second});
    \item $C_{rr}$ é o coeficiente de resistência ao rolamento;
    \item $\theta$ é o ângulo de inclinação da via (derivado do gradiente
        percentual da rota GPX);
    \item $a$ é a aceleração linear;
    \item $\eta$ é a eficiência da transmissão mecânica.
\end{itemize}

O motor físico do \argus{} resolve numericamente a equação para a velocidade
dada uma potência de entrada, utilizando o método da bisseção sobre o intervalo
de velocidades possíveis. Essa abordagem evita problemas de divisão por zero e
é numericamente estável mesmo em descidas íngremes com potência nula.

\section{Métricas fisiológicas e controle de carga de treino}
\label{sec:fisiologia}

A fisiologia do esporte consolidou métricas capazes de quantificar a carga de
treinamento e de acompanhar a evolução do atleta \cite{coggan2010}. O \argus{}
calcula todas elas localmente, a partir do vetor de potência de
\SI{1}{\hertz} registrado durante e após a sessão.

A Potência Normalizada (NP) pondera o esforço percebido em treinos variáveis
por médias móveis de 30 segundos elevadas à quarta potência:

\begin{equation}\label{eq:np}
NP = \sqrt[4]{\frac{1}{N} \sum_{i=1}^{N} p_{30,i}^4}
\end{equation}

\noindent em que $p_{30,i}$ é a i-ésima média móvel de 30 segundos do vetor de
potência. O Fator de Intensidade (IF) relaciona a NP à potência funcional
limiar (FTP) do atleta, e o \textit{Training Stress Score} (TSS) sintetiza
duração e intensidade numa escala proporcional ao esforço de uma hora no
limiar:

\begin{equation}\label{eq:if}
IF = \frac{NP}{FTP}
\end{equation}

\begin{equation}\label{eq:tss}
TSS = \frac{d \cdot NP \cdot IF}{FTP \cdot 3600} \times 100
\end{equation}

\noindent em que $d$ é a duração do treino em segundos. Complementa essas
medidas o \textit{Training Impulse} (TRIMP), que pela fórmula de Banister
pondera duração e intensidade relativa pela frequência cardíaca do atleta e
dá visão também ao componente cardiovascular do esforço.

Essas cargas alimentam o \textit{Performance Management Chart} (PMC), cujo
decaimento exponencial descreve o estado de treinamento por três curvas: a
carga crônica (CTL, \textit{fitness}), a carga aguda (ATL, \textit{fadiga}) e
o balanço entre ambas (TSB, \textit{forma}) \cite{coggan2010}:

\begin{equation}\label{eq:ctl}
CTL_{n} = CTL_{n-1} + \frac{TSS_n - CTL_{n-1}}{42}
\end{equation}

\begin{equation}\label{eq:atl}
ATL_{n} = ATL_{n-1} + \frac{TSS_n - ATL_{n-1}}{7}
\end{equation}

\begin{equation}\label{eq:tsb}
TSB = CTL - ATL
\end{equation}

\noindent em que $CTL$ representa o condicionamento acumulado, $ATL$ a fadiga
recente e $TSB$ a forma atual do atleta.

Dois indicadores complementares fecham o acompanhamento. O desvio
cardiovascular (\textit{Aerobic Decoupling}, Pw:HR) compara a relação entre
potência e frequência cardíaca nas duas metades de um treino de ritmo
constante e, acima de 5\%, alerta para desidratação ou fadiga excessiva:

\begin{equation}\label{eq:decoupling}
Decoupling = \left( \frac{P_{\text{2ª metade}}}{HR_{\text{2ª metade}}} \div \frac{P_{\text{1ª metade}}}{HR_{\text{1ª metade}}} - 1 \right) \times 100\%
\end{equation}

A recuperação da frequência cardíaca (HRR1 e HRR2) mede a queda da frequência
nos primeiros um e dois minutos após o esforço, indicador clássico de
condicionamento cardiovascular.

\section{Engenharia de software: \textit{Clean Architecture} e padrões de projeto}
\label{sec:engenharia}

A qualidade estrutural do \argus{} assenta em dois pilares de engenharia de
software: a organização das dependências proposta pela \textit{Clean
Architecture} e o emprego disciplinado de padrões de projeto clássicos.

\subsection{A regra das dependências}
\label{sec:clean-architecture}

A \textit{Clean Architecture}, de Robert C. Martin, organiza um sistema em
camadas concêntricas cujas dependências apontam sempre para dentro: o núcleo
contém as regras de negócio e não conhece banco de dados, interface gráfica,
framework ou biblioteca de hardware; os detalhes de infraestrutura dependem do
núcleo, e nunca o contrário \cite{martin2017clean}. Martin chama essa inversão
de regra das dependências, e dela derivam três consequências práticas: o
domínio torna-se independente do meio de execução, os componentes externos
tornam-se substituíveis e o núcleo passa a ser testável sem nenhum artefato
de infraestrutura ao redor.

No \argus{}, essa separação é materializada por pacotes Go com
responsabilidades distintas. O pacote \code{internal/domain} concentra as
entidades do sistema (perfis, sessões, componentes da bicicleta, planos de
treino) e as interfaces que contratam o que o domínio precisa, sem declarar
nenhuma implementação. O pacote \code{internal/usecase} implementa os casos de
uso, orquestrando esses serviços a partir das interfaces. Os pacotes
\code{internal/service} e \code{internal/repository} abrigam os adaptadores
concretos: conexão Bluetooth, motor de simulação, geração de arquivos
\texttt{.FIT}, integrações externas e IA local de um lado; persistência SQLite
do outro. O \textit{frontend} acessa esse núcleo exclusivamente pelos
contratos expostos pela \textit{bridge} Wails, sem qualquer conhecimento
sobre a estrutura interna do \textit{backend}.

\subsection{Padrões de projeto: \textit{Repository}, \textit{Facade} e \textit{Observer}}
\label{sec:padroes}

\paragraph{Repository.} O padrão \textit{Repository} isola o acesso a dados
atrás de uma interface definida pelo próprio domínio: o código de negócio
solicita persistência em termos de entidades (gravar sessão, buscar histórico)
e permanece alheio ao banco concreto, ao esquema e ao SQL \cite{gof1995}. No
\argus{}, sete contratos desse tipo vivem em
\code{internal/domain/repository.go}, um para cada categoria de dado do
sistema: perfis de usuário, atividades, potência, componentes da bicicleta,
IA, eventos e conexões Bluetooth. Suas implementações SQLite residem em
\code{internal/repository/sqlite}. A consequência prática é a
substitutibilidade: um repositório pode ser trocado por uma versão em memória
nos testes e novos repositórios, como o de desgaste de componentes criado
após o registro, entraram como módulos isolados, sem alterar os casos de uso.

\paragraph{Facade.} O padrão \textit{Facade} oferece a um subsistema de muitas
operações uma interface única e reduzida \cite{gof1995}. O
\code{StorageFacade} do \argus{} reúne dezenas de operações de armazenamento
(perfis, atividades, metas, conquistas, credenciais e fila de sincronização)
por trás de uma API compacta, poupando os casos de uso de conhecer a
orquestração entre repositórios e as regras de negócio associadas ao
salvamento. A troca da base ativa ao trocar de perfil ocorre justamente nessa
fachada (\code{LoadUserDatabase}), o que mantém o esquema multiperfil
invisível ao restante da aplicação.

\paragraph{Observer e barramento de eventos.} O padrão \textit{Observer}
desacopla quem publica um evento de quem o consome \cite{gof1995}. No \argus{}
ele aparece no \code{TelemetryEventBus} do \textit{frontend}, barramento que
publica cada pacote de telemetria vindo do \textit{backend} a todos os
assinantes (HUD, gráficos e mapa) sem saber quem os consome e sem bloquear
quem produz. Essa dissociação entre a chegada dos dados e a renderização é o
que mantém a fluididez da interface mesmo com os sensores atualizando a
telemetria continuamente, e sustenta a meta de latência definida para o
sistema.

A soma desses princípios explica por que o \argus{} absorveu a evolução
descrita no capítulo 3 sem reescritas: cada funcionalidade acrescida após o
registro entrou como um novo adaptador ou um novo caso de uso, à margem dos
módulos já existentes, preservando o comportamento das partes estáveis.

\section{Inteligência artificial generativa local}
\label{sec:ia-local}

A inteligência artificial generativa, baseada em modelos de linguagem de grande
escala (LLM), tornou-se acessível para execução local em hardware de consumo
graças a ferramentas como o Ollama \cite{ollama2024}, que empacota e serve
modelos quantizados diretamente no dispositivo do usuário. Modelos de pequena
escala (na ordem de 3 bilhões de parâmetros) conseguem executar de forma
responsiva em CPUs convencionais, sem envio de dados a servidores externos.

No \argus{}, essa capacidade é explorada por um assistente de treino que
interpreta o perfil do atleta (FTP, nível, histórico recente) e gera planos de
treino estruturados em formato ZWO, preservando a privacidade dos dados
biométricos e eliminando a dependência de serviços de nuvem para
funcionalidades de \textit{coaching}.
